% !TEX root = ../main.tex

\chapter{Introduction}
\label{ch:introduction}

Speaker verification -- determining whether two recordings were spoken by
the same person -- has become a standard component of voice-based
authentication systems. Modern systems based on deep speaker embeddings,
such as x-vectors, achieve low error rates when training and test audio
come from matched recording conditions \citep{snyder2018}.

\section{Motivation}
In practice, enrollment and verification audio are frequently captured on
different devices or transmitted over different channels -- a phone
enrolled over a landline, verified later over a mobile network, for
example. This channel mismatch introduces a systematic shift in the
learned speaker embedding space that is difficult to distinguish from a
genuine speaker difference \citep{kenny2010}.

\section{Contributions}
This thesis makes the following contributions:
\begin{itemize}
  \item A representation learning method that separates speaker identity
    from channel characteristics using paired multi-channel recordings of
    the same speakers.
  \item A benchmark constructed from recordings across five distinct
    channel conditions, used to evaluate verification performance under
    controlled channel mismatch.
  \item An evaluation showing substantially reduced equal error rate under
    mismatch relative to a standard x-vector baseline, with no loss of
    accuracy in the matched condition.
\end{itemize}

\section{Thesis Outline}
\ref{ch:related} reviews prior approaches to channel-robust speaker
verification. Later chapters (omitted from this starter document) would
present the proposed method, the benchmark construction, and experimental
results in detail.
